\section{Методология численного эксперимента}

\subsection{Цель эксперимента}

Эксперимент направлен на проверку гипотезы о том, что алгебраические
свойства модели (линейность LSA и нелинейность MiniLM) определяют
геометрию её векторного пространства и, как следствие, сильные и
слабые стороны на различных задачах.
Проверка осуществляется посредством интринсивных и экстринсивных метрик,
описанных в разделах~\ref{ssec:intrinsic} и~\ref{ssec:extrinsic}.

% ---------------------------------------------------------------
\subsection{Спецификация моделей}

\begin{table}[H]
  \centering
  \label{tab:models}
  \begin{tabularx}{\linewidth}{lX}
    \toprule
    \textbf{Компонент} & \textbf{Параметры} \\
    \midrule
    \multicolumn{2}{l}{\textit{LSA-конвейер}} \\
    \quad TF-IDF       &
    \makecell[tl]{\texttt{sublinear\_tf=True}, \texttt{min\_df=2},\\
    \texttt{max\_df=0.95}} \\
    \quad TruncatedSVD &
    $k \in \{50,\,100,\,200,\,300,\,500\}$,\ \texttt{random\_state=42} \\
    \quad Нормализация & L2 \\
    \midrule
    \multicolumn{2}{l}{\textit{all-MiniLM-L6-v2}} \\
    \quad Архитектура  & 6 слоёв, $d = 384$ \\
    \quad Агрегация    & Mean pooling с маской внимания \\
    \quad Нормализация & L2 \\
    \bottomrule
  \end{tabularx}
  \caption{Спецификация моделей}
\end{table}

% ---------------------------------------------------------------
\subsection{Датасеты}

\textbf{STS Benchmark (STSb)}~--- стандартный корпус из 8\,628 пар предложений
с оценками семантической близости от 0 до 5~.
Используется тестовая часть (1\,379 пар).

\textbf{20 Newsgroups}~--- коллекция новостных постов из 20 тематических
категорий~. Используются все 20 категорий. Набор содержит как
тематически близкие группы (\texttt{sci.*}, \texttt{rec.sport.*}),
так и контрастные (\texttt{alt.atheism}\,/\,\texttt{rec.sport.hockey}),
что позволяет проверить, насколько хорошо модель разделяет как
далёкие, так и смежные темы.

\textbf{MRPC}~--- корпус из 5\,801 пары предложений с бинарными метками
парафраз (3\,900~--- парафразы, 1\,901~--- не парафразы)~.

% ---------------------------------------------------------------
\subsection{Интринсивные метрики}
\label{ssec:intrinsic}

Данные метрики характеризуют геометрию пространства эмбеддингов независимо от
конкретной прикладной задачи.

\subsubsection{Participation Ratio}

\begin{equation*}
  \mathrm{PR}(\mathbf{E}) =
  \frac{\bigl(\sum_{i} \lambda_i\bigr)^2}{\sum_{i} \lambda_i^2},
\end{equation*}
где $\lambda_i$~--- собственные значения матрицы ковариаций
$\mathbf{C} = \tfrac{1}{n}\mathbf{E}^{\top}\mathbf{E}$ (после вычитания
среднего), $\mathbf{E} \in \mathbb{R}^{n \times d}$~--- матрица эмбеддингов.
Высокий PR означает, что дисперсия распределена по многим направлениям, а не
сосредоточена в нескольких главных компонентах. PR чувствителен прежде
всего к крупным компонентам спектра.

\subsubsection{Effective Rank}

\begin{equation*}
  \mathrm{ER}(\mathbf{E}) =
  \exp\!\Bigl(-\sum_{i} p_i \log p_i\Bigr),
  \qquad p_i = \frac{\sigma_i}{\sum_j \sigma_j},
\end{equation*}
где $\sigma_i$~--- сингулярные числа $\mathbf{E}$, $p_i$~--- нормированные
веса, задающие дискретное распределение. ER есть экспонента энтропии Шеннона
спектра сингулярных чисел; достигает максимума при равномерном распределении.
В отличие от PR, ER учитывает весь спектр, включая малые компоненты,
и поэтому обычно даёт более высокие значения.

\subsubsection{Isotropy Score}

\begin{equation*}
  \mathrm{IS}(\mathbf{E}) = \frac{1}{|I|}
  \sum_{\omega \in I}
  \frac{
    \min_{\mathbf{v}} \exp\!\bigl(f(\omega)^{\top}\mathbf{v}\bigr)
  }{
    \max_{\mathbf{v}} \exp\!\bigl(f(\omega)^{\top}\mathbf{v}\bigr)
  }.
\end{equation*}
IS~$= 1$ при идеальной изотропии (равномерном распределении на гиперсфере),
IS~$\to 0$ при концентрации в низкоразмерном подпространстве.
В вычислительной реализации используется аппроксимация
$\mathrm{IS} \approx \lambda_{\min} / \lambda_{\max}$.
IS~--- наиболее жёсткая из трёх метрик: достаточно одного слабого
направления, чтобы IS оказался близок к нулю, даже если остальные
направления используются равномерно.

\subsubsection{Hubness Score}

\begin{equation*}
  \mathrm{HS} = \frac{S_{N_k}}{\mu_{N_k}},
\end{equation*}
где $S_{N_k}$~--- стандартное отклонение, $\mu_{N_k}$~--- среднее числа
вхождений точки в $k$-NN-списки. Значение HS~$\gg 1$ указывает на
концентрацию структуры \enquote{хаббов}: точек в пространстве,
которые оказывается ближайшим соседом для непропорционально большого
числа других точек.

% ---------------------------------------------------------------
\subsection{Экстринсивные метрики}
\label{ssec:extrinsic}

\subsubsection{Семантическое сходство (STS)}

Мерой качества служит корреляция Спирмена $\rho$ между косинусным сходством
эмбеддингов и человеческими оценками. Для выборки из $n$ пар без связанных
рангов коэффициент вычисляется как

\begin{equation*}
  \rho = 1 - \frac{6\sum_{i=1}^{n} d_i^2}{n(n^2 - 1)},
\end{equation*}
где $d_i$~--- разность рангов $i$-й пары.

\subsubsection{Классификация}

В качестве классификатора используется логистическая регрессия
(\texttt{LogisticRegression}, \texttt{C=1.0}, \texttt{solver=lbfgs}).
Основной метрикой является доля верных предсказаний (Accuracy). Для учёта
возможного дисбаланса классов дополнительно вычисляется макроусреднённый
$F_1$-score:

\begin{equation*}
  \label{eq:f1}
  F_1 = 2 \cdot \frac{P \cdot R}{P + R},
  \qquad
  F_{1}^{\mathrm{macro}} = \frac{1}{C}\sum_{c=1}^{C} F_{1,c},
\end{equation*}
где $P$ и $R$~--- точность и полнота для каждого из $C$ классов.
Такое усреднение позволяет оценить качество модели на каждой тематике
равновесно.

\subsubsection{Кластеризация}

Применяется алгоритм $k$-means ($k = 10$, 10 перезапусков).
Качество оценивается двумя метриками.

\textbf{Silhouette Score.} Для каждой точки $i$ вычисляются среднее
расстояние до точек своего кластера~$a(i)$ и до точек ближайшего соседнего
кластера~$b(i)$:
\begin{equation*}
  \label{eq:silhouette}
  s(i) = \frac{b(i) - a(i)}{\max\{a(i),\, b(i)\}}.
\end{equation*}
Среднее значение по всем документам лежит в $[-1,\,1]$; значения выше $0{,}5$
свидетельствуют о выраженной кластерной структуре.

\textbf{Adjusted Rand Index (ARI).} Оценивает согласованность полученного
разбиения с истинными тематическими метками с поправкой на случайное
совпадение. Значение $\mathrm{ARI} = 1$ соответствует идеальному совпадению,
$\mathrm{ARI} \approx 0$~--- случайному разбиению.

\subsubsection{Обнаружение парафраз (MRPC)}

Порог косинусного сходства $\tau$ используется как бинарный классификатор.
Оптимальный $\tau^{*}$ выбирается по валидационному множеству из 101 значения
равномерной сетки на $[0,\,1]$. Итоговая метрика~--- $F_1$-score по
положительному классу (парафразы), которая устойчивее к дисбалансу классов,
чем Accuracy.
